\section{Methods}
\label{sec:methods}

\subsection{The autonomous pipeline}
\label{sec:methods-pipeline}
The pipeline maps a corpus of political documents to party-level estimates of scientific-evidence use
in three stages, with no human annotation and no paid API.

\paragraph{Claim extraction.} Each document is processed by an instruction-tuned LLM prompted to
extract up to ten empirical claims---propositions ``for which, if you were writing a scientific paper,
you would need a reference''---as structured records (document id, claim number, claim text), isolating
verifiable propositions from rhetorical or procedural text \citep{grimmer2013representational}.

\paragraph{Evidence retrieval.} For each claim we retrieve candidate scientific references from
\emph{OpenAlex} \citep{priem2022openalex}, a free and fully open index of over 250 million scholarly
works. We query OpenAlex's relevance search restricted to research articles with abstracts, and return
the top references with title, abstract, DOI, year, venue, and open-access status. Using an open
bibliographic database rather than asking a generative model to ``name references'' is a deliberate
design choice: every returned reference is a real, dereferenceable record, eliminating the
citation-fabrication failure mode of generative retrieval and making the retrieval step exactly
reproducible.

\paragraph{Support assessment.} Each (claim, reference) pair is assigned a stance---\textsc{support},
\textsc{refute}, or \textsc{silent} (irrelevant or insufficient information)---by the validated judges
described next. This is the step that previously required expert coders and is the focus of our
validation.

\subsection{Validators}
\label{sec:methods-validators}
We treat stance assignment as an annotation problem and use validators that fail in different ways, so
that their agreement is informative \citep{zheng2023judging}.
\begin{itemize}
  \item \textbf{LLM judge.} A high-capacity instruction-tuned model (Claude Opus~4.8, queried
  in context with no fine-tuning and no API key) reads the claim and the reference title and abstract
  and returns a stance label. Models of this class match or exceed trained human coders on stance tasks
  \citep{gilardi2023chatgpt,heseltine2024llm,tornberg2025llm}.
  \item \textbf{Open-source LLM panel.} Three instruction-tuned models from three independent
  developers---Qwen2.5-3B-Instruct (Alibaba), Phi-3.5-mini-instruct (Microsoft), and
  OLMo-2-7B-Instruct (Allen Institute for AI)---run the identical stance prompt locally. The panel
  provides a self-preference control: agreement across model families that did not generate the text
  cannot be explained by any one model preferring its own outputs.
  \item \textbf{NLI classifiers.} Two open-source natural-language-inference models with different
  architectures and training data---\texttt{DeBERTa-v3-base-mnli-fever-anli} and
  \texttt{bart-large-mnli}---score each pair with the reference as premise and the claim as hypothesis,
  mapping \textsc{entailment}$\rightarrow$\textsc{support},
  \textsc{contradiction}$\rightarrow$\textsc{refute}, \textsc{neutral}$\rightarrow$\textsc{silent}
  \citep{yin2019benchmarking,laurer2024less}. NLI is the standard discriminative approach to scientific
  claim verification \citep{wadden2020fact,koneru2024llm} and supplies a cheap surrogate that scales to
  every pair.
  \item \textbf{Baselines.} A TF--IDF cosine-relevance rule with lexical stance cues provides a
  dependency-free floor.
\end{itemize}
All models run locally on a single consumer GPU (Apple M-series, Metal backend); the heaviest component
is a 7B-parameter model in half precision.

\subsection{Human-anchored validation}
\label{sec:methods-validation}
The central methodological move is to validate every judge against expert-human labels for the
\emph{identical} task before trusting it on the unlabeled policy corpus. We use two benchmarks. SciFact
\citep{wadden2020fact} comprises expert-written claims paired with cited abstracts and human
\textsc{support}/\textsc{contradict}/\textsc{noinfo} labels; we form one (claim, abstract) pair per
cited document, yielding 340 labeled pairs. Climate-FEVER \citep{diggelmann2020climate} comprises
real-world, often contested climate claims with retrieved evidence sentences, each human-labeled
\textsc{supports}/\textsc{refutes}/\textsc{not-enough-info}; we evaluate on a class-balanced sample.
Both label schemes map onto our \textsc{support}/\textsc{refute}/\textsc{silent} classes.

We run the NLI classifiers, the open-source panel, and the TF--IDF baseline on every benchmark pair
programmatically. For the high-capacity judge, whose use we wish to keep deliberately expensive and
sparse (mirroring how it is used as gold-on-a-sample in the pipeline), we draw a class-stratified
sample, present each (claim, reference) pair \emph{blind to the gold label}, and record the judge's
stance. We then report each validator's accuracy, macro-averaged $F_1$, per-class $F_1$, and confusion
matrix against the human gold, together with pairwise inter-validator agreement (Cohen's $\kappa$). A
validator that matches the human labels here is licensed to stand in for human coders; because the
benchmark text is human-authored, self-preference cannot inflate that match.

\subsection{Valid estimation under imperfect labels}
\label{sec:methods-ppi}
Plugging machine labels directly into a downstream analysis biases estimates and invalidates confidence
intervals \citep{egami2023dsl}. We therefore treat the scalable NLI labels as a \emph{surrogate}
available for all $N$ pairs and the LLM-judge labels as \emph{gold} on a sample of $n \ll N$ pairs, and
combine them with prediction-powered inference \citep{angelopoulos2023ppi}. For the binary support
indicator, the PPI estimate of the population support rate is
\begin{equation}
  \hat{\theta}_{\text{PPI}} = \underbrace{\frac{1}{N}\sum_{i=1}^{N} f(x_i)}_{\text{surrogate mean (all pairs)}}
  + \underbrace{\frac{1}{n}\sum_{j=1}^{n}\bigl(y_j - f(x_j)\bigr)}_{\text{rectifier (gold sample)}},
  \label{eq:ppi}
\end{equation}
with $f(x_i)\in\{0,1\}$ the NLI ``support'' prediction, $y_j\in\{0,1\}$ the LLM-judge gold, variance
$\widehat{\mathrm{Var}}=\mathrm{Var}(f(x))/N + \mathrm{Var}(y-f(x))/n$, and a Wald 95\% interval
$\hat{\theta}_{\text{PPI}}\pm 1.96\,\widehat{\mathrm{Var}}^{1/2}$. Equation~\eqref{eq:ppi} is the
prediction-powered special case of the design-based supervised learning estimator \citep{egami2023dsl};
both reduce to the human-coded estimate as $n\to N$. Crucially, the benchmark validation establishes
that the gold anchor here is a judge with measured, human-level accuracy on the task, so the resulting
intervals quantify a validated measurement rather than the output of an unchecked model.

\subsection{Corpus and reproducibility}
\label{sec:methods-corpus}
The political application uses a corpus of California State Assembly member press releases (2015--2025),
partitioned by the issuing member's party, of which the AI-related subset is analysed here; a parallel
K--12 education-policy corpus is used to test domain transfer. All code, the benchmark harness, the
LLM-judge labels, the retrieved references, and the figure scripts are released at
\url{https://github.com/bdesmarais/ai-policy-science-use}; every number is reproducible from committed
outputs with open models and no paid API.
